\documentclass[11pt,a4paper,fontset=fandol]{ctexart}

\usepackage[a4paper,top=2.35cm,bottom=2.55cm,left=2.3cm,right=2.3cm,headheight=18pt,headsep=0.55cm,footskip=26pt]{geometry}
\usepackage{amsmath,amssymb,amsthm}
\usepackage{fancyhdr}
\usepackage{lastpage}
\usepackage{enumitem}
\usepackage{titlesec}
\usepackage{xcolor}
\usepackage{hyperref}
\usepackage{bookmark}
\usepackage[most]{tcolorbox}
\usepackage{setspace}
\usepackage{array}

\hypersetup{
  colorlinks=true,
  linkcolor=blue!50!black,
  urlcolor=blue!50!black,
  citecolor=blue!50!black
}

\setstretch{1.15}
\setlength{\parindent}{2em}
\setlength{\parskip}{0.28em}
\allowdisplaybreaks
\setlist[enumerate]{itemsep=0.35em, topsep=0.35em, leftmargin=2.2em}
\setlist[itemize]{itemsep=0.25em, topsep=0.25em, leftmargin=1.9em}

\definecolor{mainblue}{RGB}{36,83,140}
\definecolor{lightblue}{RGB}{241,246,252}
\definecolor{lightgray}{RGB}{246,246,246}
\definecolor{accent}{RGB}{153,76,0}
\definecolor{rulegray}{RGB}{110,118,128}

\titleformat{\section}
  {\Large\bfseries\color{mainblue}}
  {\thesection.}
  {0.45em}
  {}
  [\vspace{0.25em}\color{rulegray}\titlerule]
\titlespacing*{\section}{0pt}{1.2em}{0.8em}
\titleformat{\subsection}{\large\bfseries\color{mainblue}}{\thesubsection}{0.5em}{}

\pagestyle{fancy}
\fancyhf{}
\fancyhead[L]{\small 错位排列提高训练题单}
\fancyhead[C]{\small 组合专题：错位排列}
\fancyhead[R]{\small 刘文博}
\fancyfoot[C]{\small 第 \thepage 页 / 共 \pageref{LastPage} 页}
\renewcommand{\headrulewidth}{0.55pt}
\renewcommand{\footrulewidth}{0.35pt}

\newtcolorbox{goalbox}[1]{
  breakable,
  colback=lightblue,
  colframe=mainblue,
  boxrule=0.7pt,
  arc=1mm,
  title=\textbf{#1},
  fonttitle=\bfseries
}

\newtcolorbox{notebox}[1]{
  breakable,
  colback=lightgray,
  colframe=black!50,
  boxrule=0.55pt,
  arc=1mm,
  title=\textbf{#1},
  fonttitle=\bfseries
}

\newtheoremstyle{formalexample}
  {0.8em}{0.8em}{\normalfont}{}{\bfseries\color{accent}}{．}{0.6em}{}
\theoremstyle{formalexample}
\newtheorem{exercise}{题目}[section]
\newenvironment{solution}{\par\noindent\textbf{解：}}{\hfill$\square$\par}

\newcommand{\Dn}{D_n}

\begin{document}

\thispagestyle{empty}
\begin{titlepage}
\centering
\vspace*{0.9cm}
\begin{tcolorbox}[
  enhanced,
  width=0.92\textwidth,
  colback=white,
  colframe=mainblue,
  boxrule=1pt,
  arc=0mm,
  left=6mm,
  right=6mm,
  top=8mm,
  bottom=8mm
]
\centering
{\fontsize{24}{30}\selectfont\bfseries 错位排列提高训练题单\par}
\vspace{0.6cm}
{\fontsize{17}{22}\selectfont 适合小学高年级至初中低年级竞赛过渡训练\par}

\vspace{1cm}
\begin{tcolorbox}[colback=lightblue,colframe=mainblue,width=0.84\textwidth,boxrule=1pt]
\centering
{\large 训练目标：把“会用公式”提升到“会识别题型、会处理变式”}\\[0.35em]
{\normalsize 建议分 2--3 次完成，先独立做题，再对照详细解答订正}
\end{tcolorbox}

\vspace{1cm}
\renewcommand{\arraystretch}{1.42}
\begin{tabular}{>{\bfseries}p{3.5cm}p{8cm}}
题目数量： & 18 题，按难度递进 \\
专题重点： & 标准错位排列、恰好若干个在原位、至少若干个在原位、部分位置受限 \\
答案形式： & 末尾附较详细解答，强调思路与转化方法 \\
编译方式： & XeLaTeX \\
\end{tabular}

\vspace{0.9cm}
{\large \today\par}
\end{tcolorbox}
\end{titlepage}

\tableofcontents
\newpage

\section{使用建议}

\begin{goalbox}{做题时优先问自己的 3 个问题}
\begin{itemize}
  \item 这是标准错位排列，还是只有部分对象不能回原位？
  \item 题目问的是“全都不在原位”“恰好有几个在原位”，还是“至少有几个在原位”？
  \item 更适合用递推、公式，还是直接用
  \[
  \binom{n}{k}D_{n-k}
  \]
  这种转化？
\end{itemize}
\end{goalbox}

\begin{notebox}{建议计时}
\begin{itemize}
  \item 基础题：每题 4--6 分钟；
  \item 变式题：每题 8--12 分钟；
  \item 压轴题：先独立思考 10 分钟，再看解答。
\end{itemize}
\end{notebox}

\section{基础巩固}

\begin{exercise}
已知
\[
D_1=0,\qquad D_2=1,
\]
用递推公式求 $D_3,D_4,D_5,D_6$。
\end{exercise}

\begin{exercise}
6 个人分别拿回自己的帽子，要求没有人拿到自己的帽子，共有多少种拿法？
\end{exercise}

\begin{exercise}
数字 $1,2,3,4,5,6,7$ 排成一行，要求数字 $i$ 不在第 $i$ 个位置上，共有多少种排法？
\end{exercise}

\begin{exercise}
5 封不同的信分别装入写着对应姓名的 5 个信封里，要求每封信都不能装进自己的信封，共有多少种装法？
\end{exercise}

\section{恰好若干个在原位}

\begin{exercise}
6 个不同数字排成一行，恰好有 2 个数字在原位，共有多少种排法？
\end{exercise}

\begin{exercise}
7 封不同的信放入 7 个对应信封中，恰好有 3 封装对，共有多少种装法？
\end{exercise}

\begin{exercise}
4 个不同数字排成一行，恰好有 1 个数字在原位，共有多少种排法？
\end{exercise}

\begin{exercise}
7 个人排座位，恰好有 5 个人坐回自己的座位，共有多少种排法？
\end{exercise}

\section{至少若干个在原位}

\begin{exercise}
5 个人随机拿帽子，至少有 1 个人拿到自己帽子的拿法有多少种？
\end{exercise}

\begin{exercise}
5 封不同的信放进 5 个信封里，至少有 2 封装对，共有多少种装法？
\end{exercise}

\begin{exercise}
6 个不同数字排成一行，至多有 1 个数字在原位，共有多少种排法？
\end{exercise}

\begin{exercise}
证明：$n$ 个对象的排列中，“恰好有 $n-1$ 个对象在原位”的情况不可能发生。
\end{exercise}

\section{部分位置受限}

\begin{exercise}
4 个人甲、乙、丙、丁排座位。要求甲不坐 1 号位，乙不坐 2 号位，丙不坐 3 号位，丁没有限制。共有多少种排法？
\end{exercise}

\begin{exercise}
6 个不同数字 $1,2,3,4,5,6$ 排成一行，要求 $1,2,3$ 都不在自己的位置上，而 $4,5,6$ 没有限制。共有多少种排法？
\end{exercise}

\begin{exercise}
7 个对象排成一行，要求前 5 个对象都不在原位，后 2 个对象没有限制。共有多少种排法？
\end{exercise}

\section{综合提高}

\begin{exercise}
用容斥公式计算 $D_5$，并与递推结果核对。
\end{exercise}

\begin{exercise}
7 个人排座位，恰好有 2 个人没有坐回自己的座位，共有多少种排法？
\end{exercise}

\begin{exercise}
8 个不同数字排成一行，恰好有 1 个数字在原位，共有多少种排法？
\end{exercise}

\newpage
\section{常见易错点总结}

\begin{goalbox}{做错位排列题时，最容易丢分的 6 个地方}
\begin{enumerate}
  \item \textbf{把“恰好有 $k$ 个在原位”误算成 $\binom{n}{k}(n-k)!$。}  
  选出在原位的对象以后，剩下的对象不是“随便排”，而是必须\textbf{全部不在原位}，所以应写成
  \[
  \binom{n}{k}D_{n-k}.
  \]

  \item \textbf{把“至少有几个在原位”只算成一类。}  
  “至少有 2 个在原位”通常要拆成“恰好有 2 个、3 个、4 个……”逐类求和；或者改用补集法。

  \item \textbf{忘记“恰好有 $n-1$ 个在原位”不可能。}  
  这是错位排列题里非常高频的陷阱：如果前 $n-1$ 个都在原位，最后 1 个也只能在原位。

  \item \textbf{把“部分对象不能在原位”也直接当作 $D_m$。}  
  如果只有前 5 个对象受限制、后 2 个没有限制，那么这不是 $D_7$，也不是 $D_5$，通常应该对“坏事件”用容斥原理。

  \item \textbf{公式里的阶乘对象数看错。}  
  例如 7 个对象中若指定 2 个已经在原位，剩下应排的是 5 个对象，所以用的是 $D_5$，不是 $D_7$，也不是 $5!$。

  \item \textbf{不会检查答案是否合理。}  
  做完以后最好快速判断：答案是否小于总排法 $n!$？是否比标准错位排列 $D_n$ 大或小得合理？这样能及时发现明显错误。
\end{enumerate}
\end{goalbox}

\begin{notebox}{一个实用检查顺序}
\begin{itemize}
  \item 先判断：标准错位、恰好若干个在原位，还是部分位置受限；
  \item 再决定：用 $\binom{n}{k}D_{n-k}$、补集法，还是容斥原理；
  \item 最后检查：有没有把“不可能情形”误算进去，例如“恰好有 $n-1$ 个在原位”。
\end{itemize}
\end{notebox}

\newpage
\appendix
\section{详细解答}

\subsection*{基础巩固}

\textbf{1}\begin{solution}
错位排列的递推公式是
\[
D_n=(n-1)(D_{n-1}+D_{n-2}).
\]

因此
\[
D_3=2(D_2+D_1)=2(1+0)=2;
\]
\[
D_4=3(D_3+D_2)=3(2+1)=9;
\]
\[
D_5=4(D_4+D_3)=4(9+2)=44;
\]
\[
D_6=5(D_5+D_4)=5(44+9)=265.
\]
\end{solution}

\textbf{2}\begin{solution}
这是标准错位排列，答案就是
\[
D_6.
\]
由上一题可知
\[
D_6=265.
\]
所以共有
\[
265
\]
种拿法。
\end{solution}

\textbf{3}\begin{solution}
这也是标准错位排列，即求
\[
D_7.
\]

由递推公式
\[
D_7=6(D_6+D_5)=6(265+44)=6\times 309=1854.
\]

因此共有
\[
1854
\]
种排法。
\end{solution}

\textbf{4}\begin{solution}
5 封不同的信全部不能装入自己的信封，就是标准错位排列，所以答案为
\[
D_5=44.
\]
\end{solution}

\subsection*{恰好若干个在原位}

\textbf{5}\begin{solution}
先选出哪 2 个数字在原位，有
\[
\binom{6}{2}=15
\]
种。

剩下 4 个数字必须全部错位排列，因此有
\[
D_4=9
\]
种。

所以总数为
\[
\binom{6}{2}D_4=15\times 9=135.
\]
\end{solution}

\textbf{6}\begin{solution}
先从 7 封信中选出恰好装对的 3 封，有
\[
\binom{7}{3}=35
\]
种。

剩下 4 封必须全部装错，所以有
\[
D_4=9
\]
种。

因此总数为
\[
\binom{7}{3}D_4=35\times 9=315.
\]
\end{solution}

\textbf{7}\begin{solution}
先选哪 1 个数字在原位，有
\[
\binom{4}{1}=4
\]
种。

剩下 3 个数字必须错位排列，有
\[
D_3=2
\]
种。

所以总数为
\[
\binom{4}{1}D_3=4\times 2=8.
\]
\end{solution}

\textbf{8}\begin{solution}
若恰好有 5 个人坐回自己的座位，那么剩下 2 个人必须都不在自己的座位上。

先选出哪 5 个人在原位，有
\[
\binom{7}{5}=21
\]
种。

剩下的 2 个人只能互换位置，因此有
\[
D_2=1
\]
种。

所以总数为
\[
\binom{7}{5}D_2=21.
\]
\end{solution}

\subsection*{至少若干个在原位}

\textbf{9}\begin{solution}
总拿法有
\[
5!=120
\]
种。

没有人拿对的拿法有
\[
D_5=44
\]
种。

所以至少有 1 个人拿对的拿法为
\[
120-44=76.
\]
\end{solution}

\textbf{10}\begin{solution}
“至少有 2 封装对”可以按“恰好有几封装对”分类。

\textbf{恰好有 2 封装对：}
\[
\binom{5}{2}D_3=10\times 2=20.
\]

\textbf{恰好有 3 封装对：}
\[
\binom{5}{3}D_2=10\times 1=10.
\]

\textbf{恰好有 4 封装对：}
不可能。因为若 4 封都装对，第 5 封也只能装对。

\textbf{恰好有 5 封装对：}
有
\[
1
\]
种。

因此至少有 2 封装对的装法总数为
\[
20+10+1=31.
\]
\end{solution}

\textbf{11}\begin{solution}
“至多有 1 个在原位”包含两类：
\begin{itemize}
  \item 没有一个在原位；
  \item 恰好有 1 个在原位。
\end{itemize}

第一类有
\[
D_6=265
\]
种。

第二类：先选哪 1 个在原位，有
\[
\binom{6}{1}=6
\]
种；剩下 5 个全部错位排列，有
\[
D_5=44
\]
种。

所以第二类有
\[
6\times 44=264
\]
种。

总数为
\[
265+264=529.
\]
\end{solution}

\textbf{12}\begin{solution}
假设恰好有 $n-1$ 个对象在原位。

那么最后剩下的那 1 个对象如果不在原位，就必须去占据别人的位置；但其他 $n-1$ 个位置都已经被各自的对象占住了，它根本无处可去。

换句话说，只要前 $n-1$ 个对象都在原位，最后 1 个对象也被迫在原位。

因此“恰好有 $n-1$ 个对象在原位”是不可能的。
\end{solution}

\subsection*{部分位置受限}

\textbf{13}\begin{solution}
设
\[
A=\text{“甲坐 1 号位”},\quad
B=\text{“乙坐 2 号位”},\quad
C=\text{“丙坐 3 号位”}.
\]

总排法有
\[
4!=24
\]
种。

有
\[
|A|=|B|=|C|=3!=6,
\]
所以
\[
|A|+|B|+|C|=18.
\]

再看交集：
\[
|A\cap B|=|A\cap C|=|B\cap C|=2!=2,
\]
三项共
\[
6.
\]

三者同时发生时：
\[
|A\cap B\cap C|=1!=1.
\]

由容斥原理，满足要求的排法为
\[
24-18+6-1=11.
\]
\end{solution}

\textbf{14}\begin{solution}
这里只要求数字 $1,2,3$ 不在原位，而 $4,5,6$ 没有限制。

设
\[
A_1=\text{“1 在第 1 位”},\quad
A_2=\text{“2 在第 2 位”},\quad
A_3=\text{“3 在第 3 位”}.
\]

总排法有
\[
6!=720
\]
种。

单个坏事件的排法数：
\[
|A_1|=|A_2|=|A_3|=5!=120.
\]
三项共
\[
3\times 120=360.
\]

两个坏事件同时发生：
\[
|A_1\cap A_2|=|A_1\cap A_3|=|A_2\cap A_3|=4!=24.
\]
三项共
\[
3\times 24=72.
\]

三个坏事件同时发生：
\[
|A_1\cap A_2\cap A_3|=3!=6.
\]

因此满足要求的排法数为
\[
720-360+72-6=426.
\]
\end{solution}

\textbf{15}\begin{solution}
要求前 5 个对象都不在原位，后 2 个对象没有限制。

这和上一题完全类似，只不过现在总对象数是 7，受限制的位置有 5 个。

设前 5 个对象中，第 $i$ 个回原位的事件为 $A_i$。那么满足要求的排列数为
\[
7!-\binom51 6!+\binom52 5!-\binom53 4!+\binom54 3!-\binom55 2!.
\]

逐项计算：
\[
7!=5040,
\]
\[
\binom51 6!=5\times 720=3600,
\]
\[
\binom52 5!=10\times 120=1200,
\]
\[
\binom53 4!=10\times 24=240,
\]
\[
\binom54 3!=5\times 6=30,
\]
\[
\binom55 2!=2.
\]

所以结果为
\[
5040-3600+1200-240+30-2=2428.
\]
\end{solution}

\subsection*{综合提高}

\textbf{16}\begin{solution}
由容斥公式，
\[
D_5=5!\left(1-\frac1{1!}+\frac1{2!}-\frac1{3!}+\frac1{4!}-\frac1{5!}\right).
\]

因为
\[
5!=120,
\]
所以
\[
D_5=120\left(1-1+\frac12-\frac16+\frac1{24}-\frac1{120}\right).
\]

括号内先通分到 120：
\[
\frac{60-20+5-1}{120}=\frac{44}{120}.
\]

因此
\[
D_5=120\times \frac{44}{120}=44.
\]

与递推公式算出的结果完全一致。
\end{solution}

\textbf{17}\begin{solution}
如果恰好有 2 个人没有坐回自己的座位，那么其余 5 个人都在原位。

而这 2 个没回原位的人只能互换座位，否则仍会有人坐回原位或者无座可坐。

所以先选出哪 2 个人没有坐回原位，有
\[
\binom{7}{2}=21
\]
种；选定后他们的坐法只有
\[
1
\]
种。

故总数为
\[
21.
\]
\end{solution}

\textbf{18}\begin{solution}
恰好有 1 个数字在原位，做法是：
\begin{itemize}
  \item 先选哪 1 个数字在原位；
  \item 剩下 7 个数字全部错位排列。
\end{itemize}

所以总数为
\[
\binom81 D_7=8\times 1854=14832.
\]
\end{solution}

\end{document}
